% !TEX root = ../algebra_02.tex

\section{Канонические гомоморфизмы. Индуцированный гомоморфизм, разложение гомоморфизма}

\begin{Definition}{Канонические гомоморфизмы}{def:canonical-homs}
	Пусть $H\le G$.
	\begin{enumerate}
		\item \textbf{Гомоморфизм включения:}
		\[
		i_H\colon H\to G, \qquad h\mapsto h.
		\]
		Он инъективен, причём
		\[
		\Im i_H = H, \qquad \Ker i_H = \{e\}.
		\]
		\item Если дополнительно $H\triangleleft G$, то \textbf{канонический эпиморфизм} (проекция) задаётся формулой
		\[
		\pi_H\colon G\to G/H, \qquad g\mapsto gH.
		\]
		Это сюръективный гомоморфизм, причём
		\[
		\Im \pi_H = G/H, \qquad \Ker \pi_H = H.
		\]
	\end{enumerate}
\end{Definition}

\begin{proof}
	Гомоморфность проекции следует из равенства
	\[
	\pi_H(g_1g_2)=g_1g_2H=(g_1H)(g_2H)=\pi_H(g_1)\pi_H(g_2).
	\]
	Кроме того,
	\[
	\pi_H(g)=eH \Longleftrightarrow gH=H \Longleftrightarrow g\in H.
	\]
	Значит, $\Ker \pi_H=H$.
\end{proof}

\begin{Theorem}{Индуцированный гомоморфизм}{induced-hom}
	Пусть дан гомоморфизм
	\[
	\varphi\colon G\to G'.
	\]
	Пусть $H\triangleleft G$ и $H\subseteq \Ker\varphi$. Пусть также $H'\le G'$ и
	$\Im\varphi\subseteq H'$. Тогда отображение
	\[
	\widetilde\varphi\colon G/H\to H', \qquad \widetilde\varphi(gH):=\varphi(g),
	\]
	корректно определено и является гомоморфизмом. Более того, это единственный гомоморфизм для которого коммутирует диаграмма
\[
\begin{CD}
	G @>{\varphi}>> G' \\
	@V{\pi_H}VV @AA{i_{H'}}A \\
	G/H @>{\widetilde{\varphi}}>> H'
\end{CD}
\]
	то есть
	\[
	\varphi = i_{H'}\circ \widetilde\varphi\circ \pi_H.
	\]
\end{Theorem}

\begin{proof}
	Сначала докажем корректность.
	Если $g_1H=g_2H$, то $g_2=g_1h$ для некоторого $h\in H$. Так как
	$H\subseteq \Ker\varphi$, имеем $\varphi(h)=e$, и потому
	\[
	\varphi(g_2)=\varphi(g_1h)=\varphi(g_1)\varphi(h)=\varphi(g_1).
	\]
	Значит, значение $\widetilde\varphi(gH)$ не зависит от выбора представителя класса.

	Так как $\Im\varphi\subseteq H'$, действительно $\widetilde\varphi(gH)=\varphi(g)\in H',$ то есть отображение $\widetilde\varphi$ принимает значения в $H'$. Проверим, что $\widetilde\varphi$ является гомоморфизмом:
	\[
	\widetilde\varphi\bigl((g_1H)(g_2H)\bigr)
	=
	\widetilde\varphi(g_1g_2H)
	=
	\varphi(g_1g_2)
	=
	\varphi(g_1)\varphi(g_2)
	=
	\widetilde\varphi(g_1H)\widetilde\varphi(g_2H).
	\]

	Теперь проверим коммутативность диаграммы. Для любого $g\in G$ имеем
	\[
	(i_{H'}\circ \widetilde\varphi\circ \pi_H)(g)
	=
	i_{H'}\bigl(\widetilde\varphi(gH)\bigr)
	=
	i_{H'}(\varphi(g))
	=
	\varphi(g).
	\]
	Следовательно, 	$\varphi = i_{H'}\circ \widetilde\varphi\circ \pi_H.$

	Докажем единственность. Пусть $\psi\colon G/H\to H'$ --- гомоморфизм такой, что $\varphi = i_{H'}\circ \psi\circ \pi_H.$

	Тогда для любого $g\in G$
	\[
	i_{H'}\bigl(\psi(gH)\bigr)
	=
	(i_{H'}\circ \psi\circ \pi_H)(g)
	=
	\varphi(g)
	=
	i_{H'}\bigl(\widetilde\varphi(gH)\bigr).
	\]
	Так как $i_{H'}\colon H'\hookrightarrow G'$ --- вложение, оно инъективно. Значит,
	\[
	\psi(gH)=\widetilde\varphi(gH)
	\]
	для всех $g\in G$, откуда $\psi=\widetilde\varphi$.
\end{proof}
